\begin{helper}
Let \(s,k,m,n\) be natural numbers with \(m\le n\).  If
\(\mathrm{RamseyProperty}(s,k,m)\) holds, then
\(\mathrm{RamseyProperty}(s,k,n)\) holds.
\end{helper}
\begin{proof}
Assume \(m\le n\) and assume \(\mathrm{RamseyProperty}(s,k,m)\).  To prove
\(\mathrm{RamseyProperty}(s,k,n)\), let \(G\) be an arbitrary simple graph on
\(\operatorname{Fin}(n)\).  Let
\(\iota:\operatorname{Fin}(m)\hookrightarrow \operatorname{Fin}(n)\) be the
standard embedding obtained by viewing an index below \(m\) as an index below
\(n\).  Define a graph \(G_{0}\) on \(\operatorname{Fin}(m)\) by pulling \(G\)
back along \(\iota\): two vertices \(x,y\in\operatorname{Fin}(m)\) are adjacent
in \(G_{0}\) exactly when \(\iota(x)\) and \(\iota(y)\) are adjacent in \(G\).

By \(\noderef{RamseyProperty}\), applied to \(G_{0}\) and the assumed
\(\mathrm{RamseyProperty}(s,k,m)\), there are two cases.  In the first case
there is a finite set \(S\subseteq\operatorname{Fin}(m)\) which is an
\(s\)-clique in \(G_{0}\).  Map \(S\) into \(\operatorname{Fin}(n)\) using
\(\iota\).  Since \(\iota\) is injective, the image has the same cardinality as
\(S\), namely \(s\).  If two distinct vertices of the image are chosen, their
preimages in \(S\) are distinct.  The preimages are adjacent in \(G_{0}\)
because \(S\) is a clique, and by the definition of \(G_{0}\) their images are
adjacent in \(G\).  Therefore the image of \(S\) is an \(s\)-clique in \(G\).

In the second case there is a finite set
\(I\subseteq\operatorname{Fin}(m)\) which is a \(k\)-clique in \(G_{0}^{c}\).
Map \(I\) into \(\operatorname{Fin}(n)\) using \(\iota\).  Again injectivity
preserves cardinality, so the image has cardinality \(k\).  If two distinct
vertices of the image are chosen, their preimages in \(I\) are distinct.  Since
\(I\) is a clique in \(G_{0}^{c}\), the preimages are not adjacent in
\(G_{0}\); by the pullback definition, their images are not adjacent in \(G\).
The images are also distinct by injectivity of \(\iota\).  Hence the image is a
\(k\)-clique in \(G^{c}\).

Thus every graph \(G\) on \(\operatorname{Fin}(n)\) contains either an
\(s\)-clique or a \(k\)-clique in its complement.  This is exactly
\(\mathrm{RamseyProperty}(s,k,n)\).
\end{proof}
